\medskip

\begin{minipage}{0.48\linewidth}
On considère le cube ABCDEFGH d'arête 1.

On appelle I le point d'intersection du plan (GBD) avec la droite (EC).

L'espace est rapporté au repère orthonormé $\left(\text{A}~;~\vect{\text{AB}},~\vect{\text{AD}},~\vect{\text{AE}}\right)$.
\end{minipage} \hfill
\begin{minipage}{0.48\linewidth}
\psset{unit=1cm,linewidth=1pt}
\begin{pspicture}(7.3,5.7)
\psframe(0.8,0.3)(4.8,4.3)
\psline(4.8,0.3)(6.3,1.3)(6.3,5.3)(4.8,4.3)
\psline(6.3,5.3)(2.3,5.3)(0.8,4.3)
\psline(4.8,0.3)(6.3,5.3)
\psline[linestyle=dashed](0.8,0.3)(2.3,1.3)(6.3,1.3)
\psline[linestyle=dashed](2.3,1.3)(2.3,5.3)
\psline[linestyle=dotted,linewidth=1.2pt](4.8,0.3)(2.3,1.3)(6.3,5.3)
\psline[linestyle=dotted,linewidth=1.2pt](0.8,4.3)(6.3,1.3)
\uput[d](0.8,0.3){A} \uput[d](4.8,0.3){B} \uput[dr](6.3,1.3){C} \uput[d](2.3,1.3){D}
\uput[u](0.8,4.3){E} \uput[u](4.8,4.3){F} \uput[u](6.3,5.3){G} \uput[u](2.3,5.3){H}
\uput[dl](4.5,2.27){I} \uput[dl](3.55,0.8){J}
\psdots(4.5,2.27)(3.55,0.8)
\end{pspicture}
\end{minipage}

\medskip

\begin{enumerate}
\item Donner dans ce repère les coordonnées des points E, C, G.
\item Déterminer une représentation paramétrique de la droite (EC).
\item Démontrer que la droite (EC) est orthogonale au plan (GBD).
\item
	\begin{enumerate}
		\item Justifier qu'une équation cartésienne du plan (GBD) est :
\[x + y - z - 1 = 0.\]
		\item Montrer que le point I a pour coordonnées $\left(\dfrac23~;~\dfrac23~;~\dfrac13\right)$.
		\item En déduire que la distance du point E au plan (GBD) est égale à $\dfrac{2\sqrt 3}{3}$.
	\end{enumerate}
\item 
	\begin{enumerate}
		\item Démontrer que le triangle BDG est équilatéral.
		\item Calculer l'aire du triangle BDG.
		
On pourra utiliser le point J, milieu du segment [BD].
	\end{enumerate}
\item Justifier que le volume du tétraèdre EGBD est égal à $\dfrac13$.

\emph{On rappelle que le volume d'un tétraèdre est donné par $V = \dfrac13 Bh$ où $B$ est 
l'aire d'une base du tétraèdre et $h$ est la hauteur relative à cette base.}
\end{enumerate}
